% ============================================================
% 纯答案册·课时07-圆的方程 —— 工具/答案抽册器.py 抽册生成件（勿手改；第三档＝纯答案单独成册）
% 源件（只读）：C:/提示词/工作区/M3-第2章量产0913/成卷/导学件/课时07-圆的方程/main.tex（md5 42d73a1c44840c3d7af81728786a86b2）
%% 口径：题面不重印；逐题＝题号(印面号同正册)＋[答案]＋[详解]，值/详解照源件逐字
%       （钉值门硬断言：册值≡正册件内值≡值台账值tex，逐字全等）。册序＝正册装配序＝印面号 1..21 升序（同序同号）；键序断言基准＝题面库冻结 manifest canonical 键序。
%% 三档导出：整体 main-true／纯题 main-false（在产）｜本册＝纯答案册（本件）
%% 双档：ansbook-true.tex＝详解印本；ansbook-false.tex＝纯值速查本（详解不印）
% 编译：xelatex <壳> ×2，cwd＝本目录；sty＝qp-m3.sty 本地挂载（md5 7c3930362be8a0a2bdf21bbf8ac16573，锁校验）
% ============================================================
\documentclass[fontset=none]{ctexart}
\usepackage{qp-m3}
% —— 印面逐字符号路由补钉（承源件；− 走 TNR、▱ NSC 子块；禁改 sty 保 md5） ——
\xeCJKDeclareCharClass{Default}{"2212}%
\setCJKfamilyfont{nscblk}[Path=C:/提示词/工作区/字替对照-0909/variantF/fonts/,BoldFont=NSC-w600.ttf]{NSC-w500.ttf}%
\xeCJKDeclareSubCJKBlock{nscblk}{"25B1}%
\newcommand{\pxparallelogram}{{\CJKfamily{nscblk}▱}}%
% —— 断行微调（承母版实测档，三零门承重；\emergencystretch 不另设） ——
\xeCJKsetup{CJKglue={\hskip 0.10em plus 0.08em minus 0.05em}}%
% —— 纯答案册全线括线（长详解可跨栏断，承练习件全线括线口径；灰底盒不可跨栏断） ——
\ansblockgrayfalse
% —— 双档开关（false 档＝纯值速查：答案印、详解不印；件面开关，不动 sty 本体） ——
\newif\ifansbookdetail
\ifdefined\ansbookdetail
  \ifnum\ansbookdetail=0 \ansbookdetailfalse\else\ansbookdetailtrue\fi
\else
  \ansbookdetailtrue
\fi
% —— 页脚件名（承源件章名；件名＝<件型>答案册） ——
\renewcommand{\qpzhangming}{第二章\quad 平面解析几何}%
\renewcommand{\qpjianming}{导学件答案册}%

\begin{document}
\zhangtitle{第二章\quad 平面解析几何}
\jietitle{2.3 圆的方程}
\par\glueguard{1}\addvspace{2pt}\noindent\biaoqian{【纯答案册】}{\fangsong 题面不重印\quad 印面号与正册同号同序}\par\addvspace{2pt}

\begin{multicols}{2}
\emergencystretch=1em

\begin{ansblock}[2章-练-课时07-T1]
% ans:2章-练-课时07-T1
\ansitem{1}{D}
\ifansbookdetail
\ansnote{详解}{\((5a+1-1)^{2}+(12a)^{2}=25a^{2}+144a^{2}=169a^{2}<1\)，即\(|a|<1/13\)，故选：D．}
\fi
\end{ansblock}

\begin{ansblock}[2章-练-课时07-T2]
% ans:2章-练-课时07-T2
\ansitem{2}{D}
\ifansbookdetail
\ansnote{详解}{\((4a-1+1)^{2}+(3a+2-2)^{2}=16a^{2}+9a^{2}=25a^{2}\leqslant 25\)，即\(a^{2}\leqslant 1\)，\(-1\leqslant a\leqslant 1\)，故选：D．}
\fi
\end{ansblock}

\begin{ansblock}[2章-练-课时07-T3]
% ans:2章-练-课时07-T3
\ansitem{3}{C}
\ifansbookdetail
\ansnote{详解}{表圆需\(m>0\)。原点在圆外\(\iff(0-3)^{2}+(0+4)^{2}=25>m\)。故\(0<m<25\)，故选：C．}
\fi
\end{ansblock}

\begin{ansblock}[2章-练-课时07-T9]
% ans:2章-练-课时07-T9
\ansitem{4}{A}
\ifansbookdetail
\ansnote{详解}{圆心\((3,-5)\)。\(d=|12+15-2|/5=25/5=5\)。圆上点到直线的最短距离\(=d-r\)（\(d>r\)时）\(=1\)，得\(r=4\)，故选：A．}
\fi
\end{ansblock}

\begin{ansblock}[2章-练-课时07-T12]
% ans:2章-练-课时07-T12
\ansitem{5}{(x−1)²+y²=9（答案不唯一）}
\ifansbookdetail
\ansnote{详解}{圆心\((a,0)\)，\(r^{2}=(1-a)^{2}+9\)。取\(a=1\)：\(r=3\)，得\((x-1)^{2}+y^{2}=9\)。（一般形式\((x-a)^{2}+y^{2}=(1-a)^{2}+9\)，\(a\in\mathbf{R}\)均合．）}
\fi
\end{ansblock}

\begin{ansblock}[2章-练-课时07-T7]
% ans:2章-练-课时07-T7
\ansitem{6}{ABC}
\ifansbookdetail
\ansnote{详解}{圆心\((2,0)\)。圆关于圆心中心对称，A对；圆心在\(x\)轴上，故关于\(y=0\)对称，B对；\(x+3y-2=0\)过\((2,0)\)（\(2+0-2=0\)成立），过圆心的直线是圆的对称轴，C对；\(x-y+2=0\)：\(2-0+2=4\neq 0\)，不过圆心，D错。故选：ABC．}
\fi
\end{ansblock}

\begin{ansblock}[2章-练-课时07-T10]
% ans:2章-练-课时07-T10
\ansitem{7}{CD}
\ifansbookdetail
\ansnote{详解}{圆心\((1,-2)\)。\(l\)平分圆\(\iff l\)过圆心。截距相等分两类：①截距均为0：\(l\)过原点，\(k=-2/1=-2\)，\(l\)：\(2x+y=0\)（D对）；②截距\(a\neq 0\)：\(x/a+y/a=1\)，代入\((1,-2)\)：\((1-2)/a=1\)，\(a=-1\)，\(l\)：\(x+y+1=0\)（C对）。检验：A：\(1+2+1=4\neq 0\)，不过圆心；B：\(1+4=5\neq 0\)，不过圆心。故选：CD．}
\fi
\end{ansblock}

\begin{ansblock}[2章-练-课时07-T13]
% ans:2章-练-课时07-T13
\ansitem{8}{64；4；[3−2√2, 3+2√2]}
\ifansbookdetail
\ansnote{详解}{\((m+2)^{2}+(n+1)^{2}=|P-(-2,\ -1)|^{2}\)。圆心\(C(2,2)\)，\(r=3\)；\(|C-(-2,\ -1)|=\sqrt{16+9}=5\)。\par\noindent 距离范围\([5-3,\ 5+3]=[2,8]\)，平方范围\([4,64]\)，最大64、最小4。\par\noindent \(\sqrt{m^{2}+n^{2}}=|PO|\)，\(|OC|=2\sqrt{2}<3\)，故范围为\([3-2\sqrt{2},\ 3+2\sqrt{2}]\)．}
\fi
\end{ansblock}

\begin{ansblock}[2章-练-课时07-T4]
% ans:2章-练-课时07-T4
\ansitem{9}{B}
\ifansbookdetail
\ansnote{详解}{\(x^{2}\)与\(y^{2}\)系数均为1。\(D^{2}+E^{2}-4F=4a^{2}+16a^{2}+40a=20a(a+2)>0\)，即\(a>0\)或\(a<-2\)，故选：B．}
\fi
\end{ansblock}

\begin{ansblock}[2章-练-课时07-T5]
% ans:2章-练-课时07-T5
\ansitem{10}{A}
\ifansbookdetail
\ansnote{详解}{表圆：\(m^{2}+(-2)^{2}-4\times 2>0\)，即\(m^{2}>4\)，\(m>2\)或\(m<-2\)。点在圆外：\(1+4+m-4+2=m+3>0\)，即\(m>-3\)。取交：\((-3,-2)\cup(2,+\infty)\)，故选：A．}
\fi
\end{ansblock}

\begin{ansblock}[2章-练-课时07-T6]
% ans:2章-练-课时07-T6
\ansitem{11}{C}
\ifansbookdetail
\ansnote{详解}{过原点\(\iff 2m^{2}-6m+4=0\)，即\(m=1\)或\(m=2\)。表圆条件：\(D^{2}+E^{2}-4F=8(m-1)^{2}-4(2m^{2}-6m+4)=8m-8>0\)，即\(m>1\)。舍\(m=1\)，得\(m=2\)，故选：C．}
\fi
\end{ansblock}

\begin{ansblock}[2章-练-课时07-T8]
% ans:2章-练-课时07-T8
\ansitem{12}{D}
\ifansbookdetail
\ansnote{详解}{配方：\((x+1)^{2}+(y-m)^{2}=m^{2}+4m+5=(m+2)^{2}+1\)。\(r\) 最小为1当\(m=-2\)，此时圆心\((-1,-2)\)。\(|OC|=\sqrt{5}>r\)，圆上点到\(O\)距离最大\(=|OC|+r=\sqrt{5}+1\)，故选：D．}
\fi
\end{ansblock}

\begin{ansblock}[2章-练-课时07-T14]
% ans:2章-练-课时07-T14
\ansitem{13}{(−2,−4)；5}
\ifansbookdetail
\ansnote{详解}{表圆需\(a^{2}=a+2\)，即\(a=-1\)或\(a=2\)。\(a=2\)：\(4x^{2}+4y^{2}+4x+8y+10=0\)，除以4配方得半径平方为负，舍。\par\noindent \(a=-1\)：\(x^{2}+y^{2}+4x+8y-5=0\)，即\((x+2)^{2}+(y+4)^{2}=25\)。圆心\((-2,-4)\)，半径5．}
\fi
\end{ansblock}

\begin{ansblock}[2章-练-课时07-T11]
% ans:2章-练-课时07-T11
\ansitem{14}{x²+y²−4x+6y+8=0}
\ifansbookdetail
\ansnote{详解}{圆心在\(AB\)的中垂线\(y=-3\)上，又在\(2x-y-7=0\)上：\(2x+3-7=0\)，\(x=2\)，圆心\((2,-3)\)。\(r^{2}=4+(-3+4)^{2}=5\)。标准方程\((x-2)^{2}+(y+3)^{2}=5\)，展开：\(x^{2}+y^{2}-4x+6y+8=0\)。（验：\(A(0,-4)\)：\(16-24+8=0\)成立；\(B(0,-2)\)：\(4-12+8=0\)成立．）}
\fi
\end{ansblock}

\begin{ansblock}[2章-练-课时07-T16]
% ans:2章-练-课时07-T16
\ansitem{15}{x²+y²+2x−4y+3=0}
\ifansbookdetail
\ansnote{详解}{圆心\((-D/2,\ -E/2)\)在\(x+y-1=0\)上：\(-D/2-E/2-1=0\)，即\(D+E=-2\)。\par\noindent \(r^{2}=(D^{2}+E^{2}-12)/4=2\)，即\(D^{2}+E^{2}=20\)。由\(E=-2-D\)：\(D^{2}+(D+2)^{2}=20\)，\(2D^{2}+4D-16=0\)，\(D^{2}+2D-8=0\)，\(D=2\)或\(D=-4\)。\par\noindent \(D=2\)：\(E=-4\)，圆心\((-1,2)\)在第二象限成立；\(D=-4\)：\(E=2\)，圆心\((2,-1)\)在第四象限，舍。故\(D=2\)，\(E=-4\)，圆的一般方程为\(x^{2}+y^{2}+2x-4y+3=0\)。（验：\(r^{2}=(4+16-12)/4=2\)成立．）}
\fi
\end{ansblock}

\begin{ansblock}[2章-练-课时07-T15]
% ans:2章-练-课时07-T15
\ansitem{16}{(1) a≠−1；(2) 恒过 (0,0) 与 (16/5, −8/5)；(3) a=1/4}
\ifansbookdetail
\ansnote{详解}{(1) \(a\neq -1\)时\(x^{2}\)、\(y^{2}\)系数同为\(1+a\neq 0\)，两边除以\((1+a)\)化一般式，\(D^{2}+E^{2}-4F=(16+64a^{2})/(1+a)^{2}>0\)恒成立，表圆；\(a=-1\)时二次项消失，方程退化为直线\(-4x-8y=0\)，不合。故\(a\neq -1\)。\par\noindent (2) 方程按\(a\)整理：\((x^{2}+y^{2}-4x)+a(x^{2}+y^{2}+8y)=0\)。两定点须对一切\(a\)成立\(\iff x^{2}+y^{2}-4x=0\)且\(x^{2}+y^{2}+8y=0\)。相减：\(4x+8y=0\)，\(x=-2y\)。\par\noindent 代入\(x^{2}+y^{2}-4x=0\)：\(5y^{2}=0\)，得\((0,0)\)；代入\(x^{2}+y^{2}+8y=0\)：\(5y^{2}+8y=0\)，\(y=0\)（已得）或\(y=-8/5\)，\(x=16/5\)。定点\((0,0)\)与\((16/5,\ -8/5)\)。（验\((16/5,\ -8/5)\)：\(x^{2}+y^{2}=320/25\)，\(4x=64/5=320/25\)，\(8y=-64/5\)，两式皆零。）\par\noindent (3) 表圆时\(r^{2}=(4+16a^{2})/(1+a)^{2}\)。令\(f(a)=(4+16a^{2})/(1+a)^{2}\)（\(a\neq -1\)），求导得\(f'(a)=(32a-8)/(1+a)^{3}\)，驻点\(a=1/4\)。符号按分支定号：①\(a>1/4\)：分子\(32a-8>0\)、分母\((1+a)^{3}>0\)，\(f'>0\)；②\(-1<a<1/4\)：分子\(<0\)、分母\(>0\)，\(f'<0\)；③\(a<-1\)：分子\(<0\)，但分母\((1+a)^{3}<0\)，\(f'>0\)——该支\(f\)单调递增且\(a\to -1^{-}\)时\(f\to+\infty\)、\(f\)恒大于\(16>16/5\)，不产生最小值。\par\noindent 故\(f\)仅在\(a=1/4\)处取得最小值（\(16/5\)）。（几何：圆族中过两定点的圆以定点弦为直径时半径最小，与上互证。）故\(a=1/4\)．}
\fi
\end{ansblock}

\begin{ansblock}[2章-导-课时07-G1]
% ans:2章-导-课时07-G1
\ansitem{17}{A}
\ifansbookdetail
\ansnote{详解}{圆心\((0,b)\)：\(x^{2}+(y-b)^{2}=1\)。代入\((1,2)\)：\(1+(2-b)^{2}=1\)，得\(b=2\)。圆的方程为\(x^{2}+(y-2)^{2}=1\)，故选：A．}
\fi
\end{ansblock}

\begin{ansblock}[2章-导-课时07-G2]
% ans:2章-导-课时07-G2
\ansitem{18}{C}
\ifansbookdetail
\ansnote{详解}{\((3-2)^{2}+(2-3)^{2}=1+1=2<4=r^{2}\)，点在圆内，故选：C．}
\fi
\end{ansblock}

\begin{ansblock}[2章-导-课时07-G3]
% ans:2章-导-课时07-G3
\ansitem{19}{D}
\ifansbookdetail
\ansnote{详解}{配方：\((x+1)^{2}+(y-2)^{2}=6+1+4=11\)。圆心\((-1,2)\)，半径\(\sqrt{11}\)，故选：D．}
\fi
\end{ansblock}

\begin{ansblock}[2章-导-课时07-G4]
% ans:2章-导-课时07-G4
\ansitem{20}{(−1, 1/2)；√7/2}
\ifansbookdetail
\ansnote{详解}{两端除以2：\(x^{2}+y^{2}+2x-y-1/2=0\)，配方\((x+1)^{2}+(y-1/2)^{2}=1/2+1+1/4=7/4\)。圆心\((-1,\ 1/2)\)，半径\(\sqrt{7}/2\)．}
\fi
\end{ansblock}

\begin{ansblock}[2章-导-课时07-G5]
% ans:2章-导-课时07-G5
\ansitem{21}{−2}
\ifansbookdetail
\ansnote{详解}{圆上任意点关于\(l\)的对称点仍在圆上\(\iff\)圆关于\(l\)对称\(\iff l\)过圆心。圆心\((-1,\ -a/2)\)代入\(l\)：\(-1+a/2+2=0\)，得\(a=-2\)。（此时\(r^{2}=1+a^{2}/4+3=5>0\)，表圆成立．）}
\fi
\end{ansblock}

% ---- 尾块（\tailfill[30mm] 定高参——钉档 §三.3：无参在平衡栏呈「内容尾随框」，贴栏底须定高参；实测无参致末页平衡 Overfull\vbox 12.99pt，定高 30mm 三零） ----
\tailfill[30mm]

\end{multicols}
\end{document}
